%!TEX root = ../main.tex
% =============================================================
%  CAPÍTULO 7 — PRESENTACIÓN DE RESULTADOS
%                                       (Entrega 7 · 5 % · 13/10)
% =============================================================

\chapter{Presentación de resultados}\label{cap:resultados}

\begin{guia}[Guía de la cátedra: qué se evalúa en este capítulo]
\begin{enumerate}
    \item \textbf{Requerimientos formales}: \gls{rf} y \gls{rnf} con ID,
          descripción, prioridad y criterio de aceptación medible,
          trazados a los indicadores del capítulo~\ref{cap:metodologia}.
    \item \textbf{Evidencia de implementación}: capturas reales del sistema
          funcionando, enlaces al repositorio (anexo~\ref{anx:repositorios})
          y al despliegue.
    \item \textbf{Resultados cuantitativos honestos}: diseño experimental,
          métricas, matrices de confusión o su equivalente en el dominio,
          incluyendo lo que salió mal y su análisis. Reportar también lo
          malo y argumentarlo vale más que esconderlo. Hiperparámetros y
          versiones en el apéndice~\ref{apx:tecnicos}.
    \item \textbf{Validación con usuarios y expertos}: idealmente
          triangulada (métricas técnicas, encuesta, entrevista) y
          traducida a decisiones de diseño concretas. Instrumentos en los
          apéndices \ref{apx:encuesta} y \ref{apx:entrevista}.
\end{enumerate}
El sistema tiene que existir y andar: la entrega final incluye el
proyecto funcional.
\end{guia}

Este capítulo presenta los resultados obtenidos hasta el momento sobre el
prototipo descripto en el capítulo~\ref{cap:marco-tecnologico}: los
requerimientos definidos, la evidencia de que el sistema funciona, los
resultados de las pruebas automatizadas y la traducción de las
entrevistas y la encuesta en decisiones de diseño concretas.

\section{Requerimientos}

\begin{table}[H]
    \centering
    \caption{Requerimientos funcionales}
    \label{tab:rf}
    \small
    \begin{tabularx}{\textwidth}{l L L l L}
        \toprule
        \tb{ID} & \tb{Nombre} & \tb{Descripción} & \tb{Prioridad} & \tb{Criterio de aceptación} \\
        \midrule
        RF-01 & Autenticación por rol & El sistema permite iniciar sesión con Laravel Sanctum y asigna permisos según el rol (cliente, salón, proveedor, admin) & Alta & Un usuario sin token recibe 401 al acceder a una ruta protegida; cada rol solo ve las rutas que le corresponden \\
        RF-02 & Búsqueda pública en el marketplace & Cualquier visitante, sin iniciar sesión, puede buscar salones y proveedores y ver su disponibilidad & Alta & La búsqueda y la ficha de un salón responden sin requerir autenticación \\
        RF-03 & Solicitud de contratación & Un cliente autenticado puede solicitar un salón o proveedor, generando el evento y la reserva correspondiente & Alta & La solicitud es rechazada sin login; con login, crea evento y reserva; reutiliza el evento si el cliente ya tiene uno propio \\
        RF-04 & Recomendación de salones y proveedores & El sistema sugiere salones y proveedores según presupuesto, capacidad/ubicación y reseñas del evento del cliente & Alta & Un tenant ocupado esa fecha no aparece recomendado; uno dentro de presupuesto y capacidad puntúa más alto \\
        RF-05 & Simulación del Gemelo Digital & Un cliente puede simular cambios de invitados, fecha o presupuesto sobre su propio evento sin alterarlo, y aplicar la simulación si decide confirmarla & Alta & Un cliente no puede simular ni aplicar sobre el evento de otro; aplicar actualiza el evento real y no se puede repetir la misma simulación \\
        RF-06 & Gestión de reservas por tenant & Cada salón o proveedor administra únicamente sus propias reservas recibidas & Media & Un tenant no puede ver ni gestionar reservas de otro tenant \\
        RF-07 & Reseñas & Un cliente puede reseñar una reserva completada una sola vez, y esa reseña impacta en el puntaje de recomendación del tenant & Media & Se rechaza reseñar una reserva no completada o reseñarla dos veces \\
        \bottomrule
    \end{tabularx}
\end{table}

\begin{table}[H]
    \centering
    \caption{Requerimientos no funcionales}
    \label{tab:rnf}
    \small
    \begin{tabularx}{\textwidth}{l L L l L}
        \toprule
        \tb{ID} & \tb{Atributo} & \tb{Descripción} & \tb{Prioridad} & \tb{Umbral medible} \\
        \midrule
        RNF-01 & Seguridad de acceso & Los datos de un cliente, salón o proveedor no deben quedar expuestos a otros tenants ni a otros clientes & Alta & 0 casos donde un usuario acceda a datos ajenos en la suite de pruebas (verificado, sección~\ref{sec:pruebas}) \\
        RNF-02 & Explicabilidad de la recomendación & Cada recomendación debe venir acompañada de un motivo legible & Media & 100\,\% de las recomendaciones devuelven un campo \texttt{motivo} no vacío \\
        RNF-03 & Rendimiento & Tiempo de respuesta de las consultas principales & Media & Todavía no medido con carga real; pendiente para cuando haya despliegue (capítulo~\ref{cap:riesgos}, R-02) \\
        RNF-04 & Usabilidad & Facilidad de uso percibida por los tres tipos de usuario & Alta & Validada de forma exploratoria por encuesta y entrevistas (sección~\ref{sec:validacion}); falta prueba de usabilidad guiada sobre el prototipo \\
        \bottomrule
    \end{tabularx}
\end{table}

\section{Evidencia de implementación}

El código fuente completo está en el repositorio del sistema (ver
anexo~\ref{anx:repositorios}). \completar{Agregar capturas reales de las
pantallas principales: búsqueda del marketplace, ficha de un salón,
recomendaciones, simulador del Gemelo Digital y panel del tenant.}

\begin{figure}[H]
    \centering
    \fbox{\parbox{0.85\textwidth}{\centering\vspace{2em}
    \completar{Captura de pantalla: resultado de una búsqueda en el
    marketplace, con salones filtrados por disponibilidad.}
    \vspace{2em}}}
    \caption{Captura del sistema en funcionamiento}
    \label{fig:captura}
\end{figure}

\section{Resultados de las pruebas}\label{sec:pruebas}

A diferencia de un modelo de \gls{ml} entrenado, el motor de
recomendación es un conjunto de reglas
ponderadas, por lo que no corresponde reportar métricas como precisión o
$F_1$ sobre un conjunto de prueba: no hay ``predicción correcta o
incorrecta'' que evaluar, sino reglas de negocio que se verifican con
pruebas automatizadas de comportamiento. Lo mismo aplica al Gemelo
Digital.

Se ejecutó la suite de pruebas de integración
(\texttt{php artisan test}) sobre los cuatro módulos nuevos:

\begin{table}[H]
    \centering
    \caption{Resultado de la suite de pruebas automatizadas}
    \label{tab:tests}
    \small
    \begin{tabular}{l r r}
        \toprule
        \tb{Módulo} & \tb{Pruebas} & \tb{Resultado} \\
        \midrule
        Autenticación y roles (\texttt{AuthTest}) & 8 & 8 aprobadas \\
        Marketplace (\texttt{MarketplaceTest}) & 9 & 9 aprobadas \\
        Motor de recomendación (\texttt{RecomendacionTest}) & 8 & 8 aprobadas \\
        Gemelo Digital (\texttt{GemeloDigitalTest}) & 9 & 9 aprobadas \\
        \midrule
        \tb{Total} & \tb{34} & \tb{34 aprobadas (63 aserciones)} \\
        \bottomrule
    \end{tabular}
\end{table}

Estas pruebas verifican, entre otras cosas, que un cliente no pueda
solicitar, simular ni reseñar sobre un evento ajeno, que un tenant
ocupado en la fecha del evento no aparezca recomendado, que el costo del
catering se reescale correctamente en la simulación, y que el aviso de
capacidad excedida del salón se dispare cuando corresponde. Que la suite
pase no significa que el sistema no tenga errores: cubre los casos que
se pensaron al escribirla, no garantiza haber pensado todos los casos
posibles.

\begin{nota}
Lo que todavía no está probado, y queda pendiente: pruebas de carga o
concurrencia (RNF-03), pruebas de usabilidad guiada con usuarios reales
sobre el prototipo (más allá de la encuesta y las entrevistas), y
pruebas end-to-end del flujo completo de pagos, ya que esa
funcionalidad todavía no está integrada.
\end{nota}

\section{Validación con usuarios y expertos}\label{sec:validacion}

Los resultados se triangularon con la encuesta a 45 potenciales usuarios
y las dos entrevistas a especialistas del sector, ya presentadas en el
capítulo~\ref{cap:marco-teorico}. La tabla siguiente traduce los
hallazgos más relevantes de esa validación en decisiones concretas sobre
el sistema.

\begin{table}[H]
    \centering
    \caption{Traducción de la retroalimentación en decisiones de diseño}
    \label{tab:feedback-decisiones}
    \small
    \begin{tabularx}{\textwidth}{L L L l}
        \toprule
        \tb{Fuente} & \tb{Hallazgo} & \tb{Decisión} & \tb{Estado} \\
        \midrule
        Encuesta (n=45) & 97,8\,\% valoró útil poder simular el evento antes de confirmarlo & Priorizar el Gemelo Digital como funcionalidad central, no como agregado & Implementado \\
        Encuesta (n=45) & 93,3\,\% valoró útil la recomendación automática de proveedores & Diseñar el motor de recomendación con motivo explicable desde la primera versión (RNF-02) & Implementado \\
        Encuesta (n=45) & Solo 2,2\,\% usa hoy un marketplace existente para organizar su evento & Reforzar en el capítulo~\ref{cap:marco-teorico} que la brecha de adopción, no la falta de plataformas, es la oportunidad del proyecto & Incorporado al análisis \\
        Entrevista, dueño de salón & Los pagos son la información más difícil de mantener actualizada & Diseñar el modelo de reservas con estado y monto desde el inicio, aunque el cobro todavía se gestione manualmente (capítulo~\ref{cap:riesgos}, R-04) & Parcial \\
        Entrevista, decoración & Sería útil un plano a escala del salón para planificar la ambientación & Evaluar para una futura versión del marketplace (no forma parte del alcance actual) & Pendiente \\
        \bottomrule
    \end{tabularx}
\end{table}

\section{Resultados honestos: lo que funcionó y lo que no}

Funcionaron según lo previsto: la separación de datos por tenant (ningún
usuario accedió a datos ajenos en las pruebas), el filtro de
disponibilidad por fecha en la recomendación y en el marketplace, y la
capacidad del Gemelo Digital de detectar tanto el exceso de capacidad del
salón como el conflicto de fecha entre tenants ya confirmados.

No salió como se esperaba, o todavía no se puso a prueba:

\begin{itemize}
    \item El motor de recomendación todavía no se validó con usuarios
          reales interactuando con resultados reales, solo con las
          pruebas automatizadas de comportamiento; su precisión
          percibida por un usuario final es, por ahora, una incógnita.
    \item Las pruebas de rendimiento (RNF-03) no se hicieron porque
          todavía no hay un entorno de producción para medir con carga
          real.
    \item La entrevista a un responsable de catering, prevista en la
          Fase~2 de la metodología, sigue pendiente.
\end{itemize}
